% Chapter 7 Introduction (untitled opening paragraphs)
% Funnel structure: dN/dS gap → simulated-sky comparison → chapter roadmap → bridge to Ch. 8

The previous chapter demonstrated that the source-count distribution of faint gamma-ray emitters can be recovered from the photon-count statistics of the \textit{Fermi}-LAT sky using simulation-based inference and deep learning. The resulting $dN/dS$ extends as $\sim S^{-2}$ down to fluxes roughly twenty times below the nominal detection threshold, capturing a population of sub-threshold point sources --- blazars, star-forming galaxies, misaligned active galactic nuclei --- that contribute to the isotropic diffuse emission but remain individually undetected. Yet the $dN/dS$ is a one-dimensional flux distribution: it tells us \emph{how many} sources exist at each flux, but it says nothing about \emph{where} any specific source is located. The challenge, then, is whether this population-level statistical knowledge can be turned into spatial information.

This chapter takes the next step by complementing the source-count distribution with positional identifications. The idea is conceptually straightforward. If the $dN/dS$ is known, one can generate ensembles of realistic synthetic skies by drawing source fluxes from the measured distribution and placing them at uniformly random positions. Each synthetic sky is a plausible realization of the sub-threshold population, sharing the same statistical flux distribution as the real sky while differing in the specific source positions. By computing a pixel-wise test statistic for both the real data and each synthetic sky and comparing the resulting TS distributions via a two-sample Kolmogorov--Smirnov test, one can identify TS thresholds down to which the simulations remain statistically compatible with the data. Pixels in the real sky whose TS exceeds this threshold are labeled \emph{firing pixels} --- candidate source directions that likely host sub-threshold emitters.

Section~\ref{sec:threshold_limits} examines why the standard fixed-threshold cataloging paradigm discards useful information and produces a spatially inconsistent detection scale. Section~\ref{sec:pop_to_spatial} describes the conceptual framework for converting the $dN/dS$ into spatial source candidates through simulated-sky comparison, and introduces the frequentist meta-parameters that control the depth and purity of the resulting catalog. Section~\ref{sec:transition_catalog} transitions to the original research paper, which applies this framework to fourteen years of \textit{Fermi}-LAT data and delivers a publicly available probabilistic catalog containing approximately fifty percent more candidate source directions than the 4FGL-DR3 catalog. This extended catalog, with its globally homogeneous test-statistic scale, is well suited for the cross-correlation analyses pursued in Chapter~\ref{ch:8}, where a larger sample of source directions --- even if probabilistically defined --- improves the statistical power of collective dark matter searches.
